% 第 45 章。数字与教学模型：checks/check_ch45.py；图：figsrc/fig_ch45.py。
% 史料：monetarism_facts（取数脚本 data/fetch/monetarism45.py）；美国的 GDP、M1 与 M2：monetarism_us_1959_1999（同一脚本）。
\chapter{货币主义}\label{ch:monetarism}

\begin{chapterguide}
\guidequestion{控制货币数量就能控制通胀吗？}
\begin{guidegoals}
\item 概括货币主义的核心命题：货币需求稳定、单一规则、长期中性；
\item 说明弗里德曼与施瓦茨（1963）的货币史证据在论证中的作用；
\item 分析撒切尔政府（1979 年起）与美联储 1979—1982 年货币目标制的实践，以及“古德哈特定律”；
\item 解释货币需求不稳定为什么让各国最终放弃货币目标制。
\end{guidegoals}
% 据: fs_press, loc_thatcher, loc_thatcher_head, fomc79_date, fomc82o_date
\guideprereq{\chref{ch:naturalrate}\enspace\nameref{ch:naturalrate}；\chref{ch:money}\enspace\nameref{ch:money}}
\end{chapterguide}

1980 年 3 月 26 日是星期三。这一天，英国财政大臣豪（Geoffrey Howe，1926—2015）在下院作预算演说\cite{hansard1980,loc45}。他宣布公布政府的中期财政与货币战略，写在预算报告的第二部分，也就是人们说的“红皮书”\cite{hansard1980}。他说，这个战略的核心是一个逐年降低的货币增长目标\cite{hansard1980}。他又说，通胀若没有货币与信贷的过度扩张来迁就，长期里就维持不下去，这正是“货币主义”在实践中的核心\cite{hansard1980}。
% 据: h80_date, h80_speaker, h80_red, h80_heart, h80_monetarism, loc_howe, loc_howe_b, loc_howe_d

目标定在英镑 M3 上。英镑 M3（sterling M3）记作 £M3，是英国当时盯的广义货币\cite{hansard1986}。豪宣布，从 1980 年 2 月中旬起的 14 个月里，它按年率增长 7\% 到 11\%\cite{hansard1980}。到 1983—1984 财政年度，目标要降到 6\% 左右，是过去一年增速的一半\cite{hansard1980}。反对党议员笑起来，他回答说，这个战略绝不能和国家计划混为一谈\cite{hansard1980}。
% 据: law86_broad, h80_range, h80_months, h80_six, h80_plan, h80_laugh

一年以后，英格兰银行在 1981 年 3 月的季刊里交了账。到 1981 年 2 月中旬，£M3 一年里增长了约 20\%\cite{boe1981}。扣除“紧身衣”取消的影响，也有约 17.5\%，远在 7\% 到 11\% 的目标之上\cite{boe1981}。“紧身衣”是一种补充特别存款办法，豪在演说里宣布它 6 月中旬到期后不再延续\cite{hansard1980}。M1 是更窄的货币，主要是现金与可以开支票的存款，同期只增长 8.5\%\cite{boe1981}。
% 据: qb81_exceeded, qb81_20, qb81_175, h80_corset, qb81_m1

物价却降了下来。英格兰银行行长利-彭伯顿（Robin Leigh-Pemberton，1927—2013）1986 年回顾说，通胀的高峰超过 20\%\cite{leighpemberton1986,loc45}。1982 年 4 月，通胀降到 10\% 以下，1986 年 3 月又降到 5\% 以下\cite{leighpemberton1986}。被当作目标的那个货币涨得远比计划快，通胀照样降了（本书的解读）。控制货币数量就能控制通胀吗？
% 据: lp86_peak, lp86_peak2, loc_leighpemberton, loc_leighpemberton_head

\section{当时的世界}\label{sec:monetarism-world}

货币在经济学里失宠过一阵。弗里德曼（Milton Friedman，1912—2006）与施瓦茨（Anna J. Schwartz，1915—2012）合写过一篇长文《货币与经济周期》\cite{friedmanschwartz1963,loc45}。文章的开头回顾了这段风气，那时两人的《美国货币史》正在付印，注明 1963 年\cite{friedmanschwartz1963}。大萧条以前，人们普遍认为经济周期是一种货币现象，费雪（Irving Fisher，1867—1947）的一篇文章称之为“美元之舞”\cite{friedmanschwartz1963,loc45}。大萧条以后，凯恩斯革命让投资成了周期的主角。萨缪尔森（Paul A. Samuelson，1915—2009）在最畅销的教科书里断言，引起收入与就业波动的重要因素是投资\cite{friedmanschwartz1963,loc45}。货币即便出场，也只是被动的角色\cite{friedmanschwartz1963}。
% 据: fs_title, fs_authors, fs_press, fs_dance, fs_samuelson, fs_passive, loc_friedman, loc_schwartz, loc_fisher, loc_samuelson

他们说，近来人们重新关心货币，更多是因为通胀，周期倒在其次\cite{friedmanschwartz1963}。宽松的货币总伴随着通胀。不多少有意地限制货币存量的增长，哪里的通胀也没有被压住过\cite{friedmanschwartz1963}。钱多则物贵，这个道理在《管子》里已经有了（\chref{ch:qinhan}），货币主义把它变成了管住货币数量的政策主张。
% 据: fs_revival

他们的证据先是一个周期性的事实。1867—1960 年，美国的货币存量在经济扩张时增长，在经济收缩时多半也增长\cite{friedmanschwartz1963}。这里的货币存量，是公众持有的现金加上商业银行的活期与定期存款\cite{friedmanschwartz1963}。重大的例外只有六次，下降最少的是 1937—1938 年，降了 2.4\%\cite{friedmanschwartz1963}。下降最多的是 1929—1933 年，降了 35.2\%，就是\chref{ch:depressionus}讲过的那次收缩（\secref{sec:depressionus-money}）\cite{friedmanschwartz1963}。他们说，六次的每一次都对应着一次按别的指标看也属严重的经济收缩\cite{friedmanschwartz1963}。
% 据: fs_rise, fs_money, fs_exceptions, fs_six, fs_major

由此他们提出两个命题\cite{friedmanschwartz1963}。一是货币存量增长率的明显变化，是货币收入增长率明显变化的必要且充分的条件。二是这一点对长期的变化成立，对与经济周期差不多长的变化也成立\cite{friedmanschwartz1963}。他们进一步猜测，长期的变化主要表现在物价上，短期的变化也能大大影响产出\cite{friedmanschwartz1963}。他们也说，货币解释把货币增长率为什么会变这个问题留着不答，历史上的原因本来就多种多样\cite{friedmanschwartz1963}。
% 据: fs_props, fs_conjecture, fs_open

五年以后，弗里德曼把这些证据变成了政策处方。\chref{ch:naturalrate}讲过他 1967 年 12 月的会长演说，演说的后半讲货币政策该怎样做\cite{friedman1968}。他说，货币当局要以自己能控制的量为准绳\cite{friedman1968}。拿利率或失业率作准绳，就像宇宙飞船对准了错误的星星，导航装置再精密也要迷路\cite{friedman1968}。能控制的准绳有三个，汇率、物价指数与某一个货币总量\cite{friedman1968}。
% 据: f68_date, f68_guide, f68_criterion, f68_star, f68_three

物价本身最重要。可是货币当局的行动影响物价，比影响货币总量更间接，时滞更长，长短与大小还随情况而变\cite{friedman1968}。直接盯物价，政策本身反倒可能因时停时走成为扰动\cite{friedman1968}。他因此认为，货币总量是眼下最好的准绳，选哪一个总量远不如选定一个重要\cite{friedman1968}。他还说，货币当局的通病是动得太晚、动得太猛\cite{friedman1968}。
% 据: f68_price, f68_link, f68_indirect, f68_stops, f68_total, f68_matters, f68_toolate

毛病出在时滞上。他说，货币当局按今天的情况决定行动，行动却要过六个、九个、十二个甚至十五个月才影响经济，于是踩刹车或踩油门总是踩得太重\cite{friedman1968}。他的处方是公开宣布，让一个选定的货币总量按稳定的速度增长\cite{friedman1968}。具体的增长率与具体的总量一样，不如有一个公布出来、人所共知的增长率重要\cite{friedman1968}。他甚至说，一个平均带来温和通胀或温和通缩的固定增长率，只要稳定，也胜过过去那种大起大落\cite{friedman1968}。做不到公开的规则，货币当局能自我克制、避免大起大落，也是一大改进\cite{friedman1968}。他提议的数字在\chref{ch:naturalrate}的结尾已经讲过。
% 据: f68_lags, f68_prescription, f68_known, f68_moderate, f68_ordinance

1970 年代，各国央行陆续公布了货币目标。\chref{ch:naturalrate}讲过，联邦储备 1974 年想把 M1 的增长控制在 6\% 以内。莫迪利亚尼（Franco Modigliani，1918—2003）与 Papademos 认为，那年的严厉紧缩出自不恰当的目标，目标是按货币总量与“有序市场”定的\cite{modigliani1975,loc45}。他们主张，货币政策应当以明确说出的产出与就业目标为准\cite{modigliani1975}。加拿大银行从 1975 年起以 M1 为明确的目标\cite{bouey1982a}。瑞士国家银行在 1974 年底宣布了 1975 年的目标\cite{rich2003}。联邦德国的联邦银行也从 1975 年起每年公布货币目标\cite{gerberding2004}。
% 据: mp_squeeze, mp_orderly, mp_explicit, bou_1975, rich_1974, gwz_1975, loc_modigliani, loc_modigliani_head, loc_modigliani_d

\section{货币主义的三个命题}\label{sec:monetarism-theory}

“货币主义”这个词，豪说它本不该变成政治上骂人的话\cite{hansard1980}。把弗里德曼与施瓦茨的证据和弗里德曼的处方合起来，可以概括成三个命题（本书的概括）。第一，货币需求是稳定的。第二，货币当局应当按一个公开的规则让货币稳定地增长。第三，货币在长期里是中性的，只决定物价，不决定产出与就业。
% 据: h80_abuse

第一个命题是经验上的。Goldfeld 1976 年总结说，至少在 1974 年以前，证据表明只要少数几个因素，基本上是收入与利率，再考虑到时滞，就足以解释季度的货币需求\cite{goldfeld1976}。第三个命题就是\chref{ch:naturalrate}的长期竖线。第二个命题要靠前两个支撑（本书的解读）。
% 据: gf_few

把三个命题连起来的，是数量方程。\chref{ch:silver}的式~\eqref{eq:silver-qty}写的是$g_P\approx g_M+g_V-g_Y$，\chref{ch:bullion}的交易方程式~\eqref{eq:bullion-exchange}是一个恒等式。本书用$m$记货币的增长率，$v$记流通速度的变动率，$\pi$记通胀率，$y$记实际产出的增长率，恒等式就是$m+v=\pi+y$。货币需求稳定，$v$就可以预测，央行定下想要的通胀，就能倒推出货币该增长多少：
\begin{keyeq}
m^{*}=\pi^{*}+y^{*}-v^{*} .\label{eq:monetarism-target}
\end{keyeq}
这里$\pi^{*}$是央行能容忍的通胀率，$y^{*}$是潜在产出的增长率，$v^{*}$是流通速度的趋势变动率（本书的推导）。联邦银行的货币目标正是这样算的。按 Gerberding、Worms 与 Seitz 的叙述，它依据三个基准数字，中期的价格标准、生产潜力的增长率与流通速度的趋势下降率\cite{gerberding2004}。
% 据: gwz_three

货币要控制得住，还有一个条件。央行直接管的是基础货币与准备金，货币存量还要经过式~\eqref{eq:money-mult}的乘数（\chref{ch:money}）。\chref{ch:money}讲过，在准备金稀缺、不付利息的年代，乘数变得慢得多，货币主义者依据的就是那时的经验。乘数稳定，管住准备金就大体管住了货币（本书的解读）。2008 年以后，美联储处在地板体系，准备金有利息，乘数就不再是一个行为常数了（\chref{ch:money}）。

英国也讨论过直接管基础货币。豪在 1980 年的演说里提到，他上个星期向下院提交了一份关于基础货币控制的绿皮书\cite{hansard1980}。这份绿皮书要为公众讨论怎样改进短期的货币控制提供依据\cite{hansard1980}。他与英格兰银行行长想先听取各方意见，再决定要不要作进一步的改变\cite{hansard1980}。
% 据: h80_green, h80_views

\begin{example}{倒推一个货币目标}\label{ex:monetarism-target}
设央行能容忍的通胀是 2\%，潜在产出每年增长 2.5\%，流通速度每年下降 1\%（教学用的假设）。（1）按式~\eqref{eq:monetarism-target}，货币目标是多少？（2）若流通速度实际每年下降 4\%，货币照目标增长，通胀是多少？（3）要得到 2\% 的通胀，货币该增长多少？
\solution
（1）$m^{*}=2+2.5-(-1)=5.5\%$。（2）$\pi=5.5+(-4)-2.5=-1\%$，物价反而下降。（3）$m=2+2.5+4=8.5\%$，比原来的目标高 3 个百分点（本书的计算）。流通速度下降得比预想的快，货币超出目标反倒是维持物价所必需的。
\end{example}

\begin{pitfall}[货币增长超过目标，就是货币政策放松了]
不一定。利-彭伯顿 1986 年说，货币规则若太简单，就得常常调整或推翻\cite{leighpemberton1986}。这时目标的每一次超出，哪怕金融行为的变化足以说明它，也会被看问题过于简单的人误认为政策放松\cite{leighpemberton1986}。按他的说法，1980 年以来英国广义货币的流通速度每年下降 4\%\cite{leighpemberton1986}。劳森 1986 年也说，整个 1980 年代广义货币的增长远快于货币 GDP，却没有对通胀构成威胁\cite{hansard1986}。\exref{ex:monetarism-target}的第（2）（3）问说的就是这种情形（本书的解读）。
% 据: lp86_facile, lp86_relax, lp86_velocity, law86_faster, law86_threat
\end{pitfall}

\section{盯货币还是盯利率}\label{sec:monetarism-poole}

1972 年，联邦储备理事会的普尔（William Poole，1937 年生）与 Lieberman 写道，美国的决策者已经大大转向把货币存量当作主要工具\cite{poolelieberman1972,loc45}。1970 年以前，他们控制利率，看着货币存量。1970 年以后，他们控制货币存量，看着利率\cite{poolelieberman1972}。两人说，控制利率在技术上比较容易，控制货币存量却不然，数据来得晚，还常有大的修订\cite{poolelieberman1972}。月度的货币水平能控制到 1\% 以内的标准误差，三个月的增长率按年率算，误差却可能达到 5.6\%\cite{poolelieberman1972}。
% 据: pl_authors, pl_moved, pl_before, pl_easy, pl_lag, pl_monthly, pl_three, loc_poole

对这种局面，两人列出三种回应\cite{poolelieberman1972}。第一种认为短期的波动无关紧要，只要尽量贴着目标路径，一两年里实际的增长率会很接近目标。第二种主张放弃控制货币存量，改盯真正控制得了的变量，比如基础货币、非借入准备金或某个利率。第三种是找更好的控制办法，缩小实际与目标的偏差，他们的文章做的就是这一种\cite{poolelieberman1972}。
% 据: pl_responses, pl_r1, pl_r1b, pl_r2, pl_r3

两人说，主张把货币控制得更准，并不等于信奉货币主义，正如主张改进政府的预算程序并不等于信奉财政主义\cite{poolelieberman1972}。实际上，越看重控制的技术难题的人，往往越不看重货币政策，越看重利率的稳定\cite{poolelieberman1972}。他们的理由是，货币控制不精确，会放大收入与利率两方面的波动\cite{poolelieberman1972}。不精确的代价大到一定程度，最优的利率政策甚至可能胜过最优的货币存量政策\cite{poolelieberman1972}。
% 据: pl_doctrine, pl_camps, pl_magnify, pl_superior

那么该盯哪一个？普尔 1970 年的回答，两人在这篇文章里复述了\cite{poolelieberman1972}。扰动若主要来自 IS 曲线代表的实体部门，货币存量是恰当的工具。扰动若来自 LM 曲线代表的货币部门，利率才是恰当的工具。一般地说，工具的选择取决于实际扰动与货币扰动谁更重要\cite{poolelieberman1972}。
% 据: pl_qje, pl_qje2, pl_poole0, pl_poole, pl_general

写成\chref{ch:islm}的 IS–LM。货币需求是$L=b_0+b_1Y+b_2r+v$，IS 曲线是$Y=a_0+a_1r+u$，其中$b_1>0$，$a_1<0$，$b_2<0$，$v$与$u$是随机扰动\cite{poolelieberman1972}。本书设两者互不相关。盯利率时，$r$固定，产出只受$u$的影响。盯货币时，货币供给固定，令货币需求等于供给，联立两式解出产出偏离的方差（本书的推导）：
\begin{keyeq}
\sigma^{2}_{Y,r}=\sigma^{2}_{u},\qquad
\sigma^{2}_{Y,M}=\frac{b_2^{2}\,\sigma^{2}_{u}+a_1^{2}\,\sigma^{2}_{v}}{(a_1b_1+b_2)^{2}} .\label{eq:monetarism-poole}
\end{keyeq}
普尔的损失函数在扰动互不相关时就是这个式子（本书的核对）\cite{poolelieberman1972}。货币扰动的方差$\sigma^{2}_{v}$越大，盯货币越吃亏。实体扰动的方差越大，盯利率越吃亏，因为固定的利率让需求的冲击原样传到产出上（本书的推导）。
% 据: pl_eq1, pl_v, pl_eq5, pl_u, pl_earlier, pl_loss, pl_lr

\begin{example}{两种工具}\label{ex:monetarism-poole}
设$a_1=-1$，$b_1=1$，$b_2=-0.5$，实体扰动的方差$\sigma^{2}_{u}=1$（教学用的假设）。（1）货币扰动的方差$\sigma^{2}_{v}=1$时，哪种工具让产出更稳？（2）$\sigma^{2}_{v}=4$时呢？（3）$\sigma^{2}_{v}$超过多少，盯利率就更好？
\solution
$(a_1b_1+b_2)^{2}=(-1.5)^{2}=2.25$。（1）盯货币，$\sigma^{2}_{Y,M}=(0.25+1)/2.25\approx0.56$，小于盯利率的 1。（2）$\sigma^{2}_{Y,M}=(0.25+4)/2.25\approx1.89$，大于 1，盯利率更好。（3）令$(0.25+\sigma^{2}_{v})/2.25=1$，得$\sigma^{2}_{v}=2$（本书的计算）。
\end{example}

\fig{monetarism_poole}{盯货币与盯利率时产出的方差（\exref{ex:monetarism-poole}，本书的计算）。横轴是货币扰动的方差$\sigma^{2}_{v}$，实体扰动的方差固定为 1。盯利率时产出的方差不随$\sigma^{2}_{v}$变；盯货币时随它上升，两线在$\sigma^{2}_{v}=2$处相交。参数$a_1=-1$，$b_1=1$，$b_2=-0.5$是教学用的假设。}{fig:monetarism_poole}

货币扰动变大的证据很快就来了。Goldfeld 1976 年发现，用 1974 年以前的数据拟合得很好的货币需求方程，外推到此后，持续而且明显地高估了实际的货币需求\cite{goldfeld1976}。他的基本方程动态外推十个季度，平均高估 130 亿美元（1972 年价格）\cite{goldfeld1976}。1976 年第二季度的误差达 9.8\%，按当年价格近 300 亿美元\cite{goldfeld1976}。他的文章题为《失踪的货币》\cite{goldfeld1976}。
% 据: gf_overpredict, gf_13, gf_13b, gf_98, gf_title

Goldfeld 自认，作为解释这块缺口的尝试，这篇文章“相当失败”\cite{goldfeld1976}。他的结论是，在实用的层面上，货币需求函数多半已经发生了某种移动\cite{goldfeld1976}。Goldfeld 写道，按普尔的结果，货币扰动的分量加重，政策就该向利率倾斜\cite{goldfeld1976}。他说，本文的结果明确支持眼下采用利率政策，而不是 M1 政策\cite{goldfeld1976}。
% 据: gf_failure, gf_shift, gf_poole, gf_clearly

Goldfeld 追查过的线索里，有一条与现金管理有关。主动的现金管理有固定的成本，人们还要边干边学，熟悉新的办法，交易技术于是成了持币者自己的选择\cite{goldfeld1976}。新技术一旦采用，利率即便回落，企业与个人也未必放弃，货币需求因此可能有一种棘轮效应\cite{goldfeld1976}。高利率教会了人们少持现金，利率回落以后也没有忘掉（本书的解读）。
% 据: gf_learning, gf_ratchet

\begin{compare}{货币目标与利率目标}{货币目标}{利率目标}
\cmprow{央行固定什么}{某一货币总量的增长路径}{短期利率}
\cmprow{哪种扰动下更好}{实体部门的扰动为主\cite{poolelieberman1972}}{货币部门的扰动为主\cite{poolelieberman1972}}
\cmprow{需求冲击来了}{利率自动上升，挡住一部分（本书的推导）}{冲击原样传到产出（本书的推导）}
\cmprow{货币需求移动了}{利率与产出跟着波动（本书的推导）}{货币量自动迁就，产出不受影响（本书的推导）}
\cmprow{技术上}{数据晚、常修订\cite{poolelieberman1972}}{比较容易\cite{poolelieberman1972}}
\cmpsum{两种目标的高下，取决于扰动主要来自实体部门还是货币部门。1970 年代中期以后货币需求不再稳定，天平就向利率倾斜（本书的概括）。}
\end{compare}

\section{美国：1979—1982 年}\label{sec:monetarism-us}

1979 年 10 月 6 日是星期六。联邦公开市场委员会这天开会，会是主席召集的\cite{fomc1979,loc45}。主席是沃尔克（Paul A. Volcker，1927—2019）\cite{fomc1979,loc45,gnd45}。会议要讨论怎样更好地控制货币与银行信贷的扩张，因为金融与商品市场出现了投机，通胀的力量也很强\cite{fomc1979}。委员会年初定下，1978 年第四季度到 1979 年第四季度 M1 增长 1.5\% 到 4.5\%\cite{fomc1979}。实际上，第一季度按年率下降约 2\%，第二、第三季度却按年率增长约 7.5\% 与 9.5\%\cite{fomc1979}。
% 据: fomc79_date, fomc79_called, fomc79_purpose, fomc79_purpose2, fomc79_feb, fomc79_period, fomc79_ranges, fomc79_m1q, loc_volcker_b, loc_volcker_d, loc_volcker_head, loc_volcker2, gnd_volcker_b, gnd_volcker_d, gnd_volcker_fed

大多数委员强烈支持改变公开市场操作的办法\cite{fomc1979}。新办法着重供应与目标货币增长相称的银行准备金，同时允许联邦基金利率比过去大得多地波动\cite{fomc1979}。少数委员仍想盯着联邦基金利率\cite{fomc1979}。会议记录说，主要的理由是新办法更有把握实现货币增长的目标\cite{fomc1979}。通胀很快的环境里，对利率、货币增长与经济活动之间关系的估计已经不如从前可靠\cite{fomc1979}。委员们也承认，准备金与货币总量之间的关系并不精确\cite{fomc1979}。
% 据: fomc79_shift, fomc79_few, fomc79_reason, fomc79_reliable, fomc79_precise

1979 年的转向是谁定的？委员会指示，9 月到 12 月 M1 按年率增长约 4.5\%\cite{fomc1979}。条件是联邦基金利率的周平均大体留在 11.5\% 到 15.5\% 之间\cite{fomc1979}。委员会预料，新办法会让联邦基金利率马上上升，波动也更大\cite{fomc1979}。投赞成票的是沃尔克等 12 人，没有反对票\cite{fomc1979}。会后，理事会把贴现率从 11\% 提到 12\%，对管理负债的增加额设了 8\% 的边际准备金要求\cite{fomc1979}。委员们认为，这些措施会增强国内外对联邦储备实现货币目标的决心的信心\cite{fomc1979}。
% 据: fomc79_directive, fomc79_funds, fomc79_weekly, fomc79_prompt, fomc79_vote, fomc79_after, fomc79_board, fomc79_marginal, fomc79_confidence

目标仍然难打中。委员会在 1982 年 2 月回顾 1981 年时说，按转入 NOW 账户的影响调整以后，M1-B 只增长约 2.25\%，比目标下限低约 1.25 个百分点\cite{fomc1982a}。按委员会的定义，M1 包括银行与储蓄机构的活期存款、流通中的现金、旅行支票、NOW 账户与自动转账账户，以及信用社的股金汇票账户\cite{fomc1982a}。NOW 账户是一种可以开支票的付息账户（本书的概括）。较宽的两个口径却超出了各自的上限，其中 M2 增长约 9.5\%\cite{fomc1982a}。按委员会的解释，M1-B 偏低，是因为金融创新与高利率改变了现金管理的办法\cite{fomc1982a}。
% 据: fomc82f_date, fomc82f_m1b, fomc82f_def, fomc82f_def2, fomc82f_m2m3, fomc82f_why

委员会没有放弃 M1，只是换了口径。从 1982 年起，M1-A 与按转移调整的 M1-B 不再编制\cite{fomc1982a}。留下的 M1 就是原来的 M1-B，包括全部 NOW 账户\cite{fomc1982a}。委员会重申了 1981 年年中暂定的 1982 年区间，M1 是 2.5\% 到 5.5\%\cite{fomc1982a}。委员会推测，NOW 账户最近的扩张多半是个人流动性偏好的一时反常，货币与名义国民生产总值的关系会更接近历史的样子\cite{fomc1982a}。
% 据: fomc82f_drop, fomc82f_ml, fomc82f_range82, fomc82f_aberration

1982 年 10 月 5 日，委员会又退了一步\cite{fomc1982b}。10 月初有大批全民储蓄存单（all savers certificates）到期，稍后存款机构放松管制委员会还要批准一种不受利率上限约束的新账户\cite{fomc1982b}。委员会认为，新账户对 M1 的影响多半很大，大小却无从预测\cite{fomc1982b}。它因此决定，这段时间对 M1 的变动给予远少于通常的重视，不为它定具体的增长目标\cite{fomc1982b}。联邦基金利率若持续落在 7\% 到 10.5\% 之外，主席可以召集委员们商议\cite{fomc1982b}。
% 据: fomc82o_date, fomc82o_savers, fomc82o_didc, fomc82o_basis, fomc82o_weight, fomc82o_chair, fomc82o_funds

这次表决 9 票对 3 票\cite{fomc1982b}。投反对票的 Horn（生卒年未详）主张继续为 M1 定具体的目标\cite{fomc1982b}。另外两人反对，是想在眼下更坚决地压住货币增长\cite{fomc1982b}。其中一人担心，指令可能被误读，在财政政策高度扩张的背景下有损长期的反通胀目标\cite{fomc1982b}。
% 据: fomc82o_vote, fomc82o_horn, fomc82o_dissent, fomc82o_ford, fomc82o_fiscal

代价落在谁身上？委员会 1982 年 2 月的记录写道，1981 年第四季度实际国民生产总值按年率下降约 5.25\%，制造业减少了 70 多万个工作\cite{fomc1982a}。12 月的失业率又升了 0.5 个百分点，到了 8.9\%\cite{fomc1982a}。10 月的记录写道，8 月的失业率仍是 9.8\%，9 月中旬首次申领失业保险的人数创了新高\cite{fomc1982b}。联邦储备的工作人员预测，下一年的失业仍会停在高水平\cite{fomc1982b}。这几年的利率、衰退与反通胀的总账，是\chref{ch:credibility}的题目。
% 据: fomc82f_gnp, fomc82f_jobs, fomc82f_unemp, fomc82o_aug, fomc82o_unemp, fomc82o_claims, fomc82o_staff0, fomc82o_staff

\fig{monetarism_velocity}{美国 M1 与 M2 的流通速度，1959—1999 年（本书的计算）。流通速度是当年的名义 GDP 除以 12 月的货币存量。M1 的流通速度到 1981 年大体一直上升，1982 年与 1985—1986 年急降，此后忽上忽下；M2 的流通速度直到 1990 年都在一个窄的范围里，1990 年代上升。竖线标出 1979、1982 与 1986 年。数据：《总统经济报告》（2000）表 B-1 与 B-67。}{fig:monetarism_velocity}
% 据: erp_b1_title, erp_b1_units, erp_b67_title, erp_b67_dec, erp_b67_december, erp_b67_note, erp_b1_1979, erp_b67_1979, erp_b1_1999, erp_b67_1999

流通速度的数字说明了为什么（\figref{fig:monetarism_velocity}）。按《总统经济报告》的两张表，用当年的名义 GDP 除以 12 月的 M1，可以算出流通速度\cite{erp2000b}。1959 年它约为 3.6，1981 年约为 7.2（本书的计算）。22 年里平均每年上升约 3.2\%（本书的计算）。1982 年它下降 4.3\%，1986 年又下降 9.6\%（本书的计算）。M1 的流通速度若照 1981 年以前的趋势走，同样的 M1 目标就会配上快得多的名义收入（本书的解读）。
% 据: erp_b1_1959, erp_b67_1959, erp_b1_1981, erp_b67_1981, erp_b1_1982, erp_b67_1982, erp_b1_1985, erp_b67_1985, erp_b1_1986, erp_b67_1986

\section{英国：中期财政战略}\label{sec:monetarism-uk}

英国的目标是谁定的？演说是财政大臣豪作的，他说政府继续把对付通胀当作第一要务\cite{hansard1980}。他说，以为在战胜通胀与别的路之间真有选择，是一种错觉\cite{hansard1980}。预算报告载有政府的四年计划，涵盖货币增长、公共开支与税收\cite{hansard1980}。“紧身衣”的去留，他说是与英格兰银行行长商定的\cite{hansard1980}。他预先说明，紧身衣取消以后，原来被挤到银行体系之外的信用会回流，£M3 会因此膨胀\cite{hansard1980}。这一部分的增长不代表货币状况的变化，规模也无法精确计量或预测\cite{hansard1980}。
% 据: h80_priority, h80_illusion, h80_path, h80_governor, h80_swollen, h80_signal

紧身衣为什么要取消？利-彭伯顿 1986 年解释，1979 年底英国取消了外汇管制，很快就让当时直接的国内货币管制失效，紧身衣首当其冲，第二年也取消了\cite{leighpemberton1986}。他还说，现行货币政策的基础其实是 1976 年在另一届政府手里打下的\cite{leighpemberton1986}。\chref{ch:naturalrate}讲过，同年 9 月，工党的首相在党的大会上已经改了口。
% 据: lp86_exchange, lp86_1976

结果比预想的更难看。英格兰银行 1981 年 3 月说，£M3 的年增长率秋天接近 25\%\cite{boe1981}。1980—1981 年度的公共部门借款需求（PSBR）估计是 135 亿英镑，比 1980 年 3 月的预测多 50 亿英镑\cite{boe1981}。PSBR 是公共部门一年要借的钱（本书的概括），它与货币目标的关系是\chref{ch:budget}的题目。1981 年 3 月的预算把 PSBR 定在 105 亿英镑左右\cite{boe1981}。£M3 的新目标定为 6\% 到 10\%\cite{boe1981}。英格兰银行的最低贷款利率同时从 14\% 降到 12\%\cite{boe1981}。
% 据: qb81_25, qb81_psbr, qb81_plan, qb81_target, qb81_mlr

代价落在谁身上？豪在演说里承认有人责怪政府的政策不必要地严厉，他回答说，英国的困难太深，严厉的政策无可避免\cite{hansard1980}。他说，放松货币与预算，短期里也许能多些产出，最终只会助长新的通胀，随之而来的危机会毁掉更多的工作\cite{hansard1980}。
% 据: h80_tough, h80_tough2, h80_relaxed, h80_crisis

英格兰银行在同一期季刊里写道，截至 1980 年第四季度的一年里，制造业产出下降约 15\%，是战后任何十二个月里最大的降幅\cite{boe1981}。企业在第四季度大幅削减存货与雇员，也压缩了固定投资\cite{boe1981}。按实际计算的高利率也抑制了借款\cite{boe1981}。建筑协会建议的住房抵押贷款利率先降到 14\%，后来又降到 13\%\cite{boe1981}。\chref{ch:naturalrate}讲过，英国的失业率 1981 年是 9.6\%，1983 年到了 11.5\%。
% 据: qb81_q4, qb81_war, qb81_labour, qb81_real, qb81_mortgage

1986 年 3 月 18 日，财政大臣劳森（Nigel Lawson，1932—2023）在预算演说里换了说法\cite{hansard1986,loc45,gnd45}。他说，中期财政战略的核心是逐步降低以现价 GDP 衡量的总支出的增长\cite{hansard1986}。他说，过去六年货币 GDP 的增长率减了一半\cite{hansard1986}。广义货币的目标从去年秋天起就显得定低了，新的范围提高到 11\% 到 15\%\cite{hansard1986}。他说，短期利率才是货币政策的根本工具\cite{hansard1986}。
% 据: law86_date, law86_speaker, law86_heart, law86_halved, law86_low, law86_range, law86_rates, loc_lawson_head, loc_lawson2, gnd_lawson_head, gnd_lawson_b, gnd_lawson_d

狭义货币是另一回事。劳森给狭义货币 M0 定的区间是 2\% 到 6\%，与上一年的中期财政战略预告的一样\cite{hansard1986}。他说，利率变动对狭义货币与汇率的影响相当快而直接，对广义货币的影响复杂得多，也迟缓得多\cite{hansard1986}。他又说，金融自由化与创新让货币管理更复杂了，政策的实质却没有任何改变\cite{hansard1986}。
% 据: law86_broad, law86_m0, law86_narrow, law86_essence

同年 10 月 22 日，利-彭伯顿在拉夫堡大学的演讲里讲了原因\cite{leighpemberton1986}。从二战结束到 1979 年，除了 1971—1973 年短暂而剧烈的中断，广义货币的流通速度大体一直在上升\cite{leighpemberton1986}。1980 年以后，它平均每年下降 4\%\cite{leighpemberton1986}。过去六个年度的£M3 目标，只有两个达成了\cite{leighpemberton1986}。他把货币目标比作射击，说若以为目标能精确认出、只要射手够有本事够坚决就一定打中，那是错的\cite{leighpemberton1986}。射手不知道调整瞄准，就可能重伤经济\cite{leighpemberton1986}。过去五六年，英国在最终目标上做得比在中间目标上好得多\cite{leighpemberton1986}。
% 据: lp86_date, lp86_title, lp86_lecture, lp86_issue, lp86_rose1, lp86_rose, lp86_velocity, lp86_twoofsix, lp86_marksman, lp86_adjust, lp86_success

1987 年 3 月 17 日，劳森说，中期财政战略的中心目标仍是逐步降低货币 GDP 的增长，把通胀挤出经济，最终实现物价稳定\cite{hansard1987}。短期利率仍是根本工具，要按狭义货币、广义货币与汇率所显示的货币状况来定\cite{hansard1987}。他接着宣布，正如行长在拉夫堡的演讲里有力地论证过的，眼下也许更明智的做法是完全不为广义货币设明确的目标\cite{hansard1987}。
% 据: law87_date, law87_speaker, law87_central, law87_rates, law87_eschew

加拿大走得更早。1982 年 11 月 29 日，加拿大银行行长 Bouey（Gerald K. Bouey，1920—2004）在多伦多说，M1 一年多来已经不再是可靠的指南\cite{bouey1982a,loc45}。企业的现金管理可以把活期余额隔夜投资出去，个人有了兼具支票与储蓄功能的新账户，钱于是大量而持续地流出 M1\cite{bouey1982a}。他宣布，加拿大银行不再为 M1 设目标区间，并说对此感到遗憾\cite{bouey1982a}。他同时说，通胀的下降伴随着经济活动的严重萎缩，高失业、大量闲置的设备、低利润与许多企业倒闭\cite{bouey1982a}。
% 据: bou_date, bou_reliable, bou_shifts, bou_drop, bou_regret, bou_decline, bou_strains, bou_practices, loc_bouey, loc_bouey_head

他说遗憾，是因为一个与总支出的趋势有系统联系的货币总量，有助于判断货币该扩张多少，尤其有助于避免累积的错误\cite{bouey1982a}。他又说，没有可用的货币目标，最近一段并不是要紧的事\cite{bouey1982a}。加拿大银行不肯不顾汇率，把利率一下压低几个百分点，因为那会让人对加元失去信心\cite{bouey1982a}。
% 据: bou_why, bou_critical, bou_dollar

加拿大银行一开始就知道这层风险。Bouey 在 12 月说，多年里，M1、利率与总支出的趋势之间有相当稳定的关系，加拿大银行就拿它作货币目标的根据\cite{bouey1982b}。他引了自己在 1975 年年报里的话，M1 只有在公众每一元收入所持的交易余额保持相当可预测的模式时，才能为政策提供有用的信息\cite{bouey1982b}。他当年还说，公众的这种习惯可能随着支付技术的创新而改变\cite{bouey1982b}。
% 据: boud_stable, boud_basis, boud_aware, boud_1975, boud_habits

\section{古德哈特定律}\label{sec:monetarism-goodhart}

三国的故事有一个共同点。货币总量一旦成了目标，它与收入的关系就变了（本书的概括）。1983 年，澳大利亚储备银行的 Moses 在一篇讨论稿里列出美国 1970 年代中期以来的三个问题\cite{moses1983}。货币需求的短期不稳定，M1 与经济活动的关系不再可靠，M1 的增长比它与 GDP 的历史关系所预示的慢\cite{moses1983}。第三个问题，就是 Goldfeld 说的失踪的货币（本书的解读）。Moses 在脚注里说，英格兰银行的古德哈特（Charles Goodhart，1936 年生）也许早就料到了\cite{moses1983,loc45}。
% 据: moses_title, moses_author, moses_views, moses_three, moses_goodhart, loc_goodhart, loc_goodhart_head

照古德哈特 1975 年的经验法则，任何观察到的统计规律，一旦为了控制而对它施加压力，就会趋于瓦解\cite{moses1983}。Moses 在同一篇文章里把它叫作“古德哈特定律”\cite{moses1983}。这条法则有一个近亲，就是\chref{ch:ratex}的卢卡斯批判（本书的解读）。放进普尔的式子里，它说的是货币一旦成了目标，货币扰动的方差$\sigma^{2}_{v}$就会变大，天平向利率倾斜（本书的解读）。
% 据: moses_rule, moses_ref, moses_credit

Moses 的分析用的正是 IS–LM\cite{moses1983}。按 Moses 的概括，近年有些因素改变了 LM 曲线的斜率，另一些让它移动了位置\cite{moses1983}。新的交易技术减少了对狭义货币的需求，使相应的 LM 曲线右移\cite{moses1983}。重新定义货币，也许能纠正位置的移动，前提是移动的大小可以预测，却纠正不了斜率的变化\cite{moses1983}。Moses 的结论是，金融创新使制定货币政策时更需要看一系列金融变量\cite{moses1983}。Moses 还引了明斯基 1957 年的话，金融市场的结构或运作一旦迅速变化，央行行动的效力就得重新检验\cite{moses1983}。
% 据: moses_islm, moses_slope, moses_shift, moses_redefine, moses_range, moses_minsky

上面三国的记录都能这样读（本书的解读）。豪预先知道，取消紧身衣会让£M3 膨胀\cite{hansard1980}。联邦公开市场委员会把 M1 偏低归于金融创新与高利率下的现金管理\cite{fomc1982a}。Bouey 把 M1 的失灵归于银行业务的变化\cite{bouey1982a}。劳森说广义货币增长快，是因为低通胀与高实际利率下付息存款更有吸引力，还有金融体系的创新与自由化\cite{hansard1986}。
% 据: h80_swollen, fomc82f_why, bou_shifts, law86_reasons

\section{数据检验与反例}\label{sec:monetarism-test}

如果货币目标注定失败，联邦德国与瑞士就是反例。利-彭伯顿 1986 年说，过去五年左右，美国与英国的金融变化特别快，这与它们和其他盎格鲁-撒克逊国家执行货币目标的困难也许有联系\cite{leighpemberton1986}。加拿大与澳大利亚都放弃了货币目标，美国与英国一样不得不重新定义目标\cite{leighpemberton1986}。瑞士与德国做得更成功，也许部分是因为它们由来已久的全能银行结构\cite{leighpemberton1986}。\tabref{tab:monetarism-targets}把两家央行的目标与实际放在一起。
% 据: lp86_pace, lp86_anglo, lp86_given, lp86_swiss

\begin{table}[tbp]\centering\small
\caption{联邦银行与瑞士国家银行的货币目标与实际增长（\%），1975—1998 年}\label{tab:monetarism-targets}
\begin{tabularx}{\linewidth}{@{}lYYYY@{}}
\toprule
 & \multicolumn{2}{c}{联邦银行} & \multicolumn{2}{c}{瑞士国家银行}\\
\cmidrule(lr){2-3}\cmidrule(l){4-5}
年 & 目标 & 实际 & 目标 & 实际\\
\midrule
1975 & 8 & 10 & 6 & 4.4\\
1978 & 8 & 11 & 5 & 16.2\\
1980 & 5—8 & 5 & 4 & $-0.6$\\
1983 & 4—7 & 7 & 3 & 3.6\\
1986 & 3.5—5.5 & 8 & 2 & 2.0\\
1987 & 3—6 & 8 & 2 & 2.9\\
1988 & 3—6 & 7 & 3 & $-3.9$\\
1990 & 4—6 & 6 & 2 & $-2.6$\\
1992 & 3.5—5.5 & 9 & — & —\\
1995 & 4—6 & 2 & — & —\\
1998 & 3—6 & 6 & — & —\\
\bottomrule
\end{tabularx}
\datasource{联邦银行：Gerberding、Worms 与 Seitz（2004）表 10，1987 年以前是中央银行货币存量，1988 年起是 M3；1978 年是年平均的目标与实际，其余年份是第四季度对上年第四季度（1975 年为 12 月对 12 月）。瑞士国家银行：Rich（2003）表 1，一般是月度同比增长率的年平均，1980 年相对 1979 年 11 月的水平，1990 年是第四季度水平的同比；1975—1978 年是 M1，1980—1988 年是剔除月末准备金突增的调整后基础货币，1990 年是经季节调整的基础货币；1979 年没有货币目标，表 1 只列到 1990 年。两表取自文字层，并按页面图像核对，2026 年 11 月 12 日下载。表中只列部分年份。}
% 据: gwz_t10, gwz_t10head, gwz_t10m3, gwz_t10n1, gwz_t10n2, gwz_t1975, gwz_t1976, gwz_t1977, gwz_t1978, rich_t1, rich_t1note, rich_t1m1, rich_mba, rich_mbsa, rich_mbsa_def
\end{table}

\fig{monetarism_targets}{两家央行的货币目标与实际增长（\%）。(a) 联邦银行，1979—1998 年：灰色竖条是公布的目标区间（1989 年是 5\% 的点目标），圆点是第四季度对上年第四季度的实际增长。(b) 瑞士国家银行，1975—1990 年：方块是年度目标，圆点是实际增长；1979 年没有目标。数据：Gerberding、Worms 与 Seitz（2004）表 10，Rich（2003）表 1。}{fig:monetarism_targets}
% 据: gwz_t1979, gwz_t1980, gwz_t1981, gwz_t1982, gwz_t1983, gwz_t1984, gwz_t1985, gwz_t1986, gwz_t1987, gwz_t1988, gwz_t1989, gwz_t1990, gwz_t1991, gwz_t1992, gwz_t1993, gwz_t1994, gwz_t1995, gwz_t1996, gwz_t1997, gwz_t1998, rich_t1975, rich_t1976, rich_t1977, rich_t1978, rich_t1979, rich_t1980, rich_t1981, rich_t1982, rich_t1983, rich_t1984, rich_t1985, rich_t1986, rich_t1987, rich_t1988, rich_t1989, rich_t1990

两家央行的目标由谁定？联邦银行的货币政策由中央银行理事会决定，目标通常在 12 月公布\cite{gerberding2004}。瑞士国家银行的决策由三人组成的理事会作出，理事会通常在 11 月开会，12 月由行长在记者会上宣布\cite{rich2003}。
% 据: gwz_council, gwz_december, rich_three, rich_board, rich_president, rich_decided

联邦银行的目标打得并不准（\figref{fig:monetarism_targets}a）。按 Gerberding 等人的统计，1979—1998 年的 20 年里，目标落空 7 次，其中 6 次是实际增长高于目标\cite{gerberding2004}。本书按表 10 数了一遍，落空的是 1986、1987、1988、1992、1993、1995 与 1996 年，只有 1995 年低于目标（本书的核对）。
% 据: gwz_missed, gwz_t1986, gwz_t1987, gwz_t1988, gwz_t1992, gwz_t1993, gwz_t1995, gwz_t1996

拉长了看，账却对得上。按 Gerberding 等人的计算，1979—1998 年货币的平均增长略低于 6\%，产出的趋势增长率是 2.1\%\cite{gerberding2004}。同期流通速度的趋势下降率是 1.1\%，平均通胀是 2.8\%\cite{gerberding2004}。按式~\eqref{eq:monetarism-target}倒过来算，$6-2.1-1.1=2.8$，正好对上（本书的核对）。逐年的目标常常落空，二十年的平均却符合数量方程（本书的解读）。
% 据: gwz_average

联邦银行的价格标准也一路往下调。按 Gerberding 等人的整理，1975—1998 年它从 5\% 降到 1.75\%\cite{gerberding2004}。到 1984 年为止，这个数代表联邦银行眼中当年“不可避免”的涨价幅度，1985 年起改为中期里能容忍的最大涨幅\cite{gerberding2004}。联邦银行并不惜一切去打中目标，偏差若被认为来自产出或流通速度对趋势的短期偏离，它有时就容忍了\cite{gerberding2004}。
% 据: gwz_norm, gwz_unavoidable, gwz_tolerated

\begin{case}{联邦德国与瑞士 · 1975—1999 · 守着货币目标的两家央行}
联邦银行年中复核目标，复核以后改过目标区间的只有 1991 年一次\cite{gerberding2004}。目标是 2 到 3 个百分点宽的走廊，只有 1989 年是 5\% 的点目标\cite{gerberding2004}。
% 据: gwz_midyear, gwz_corridor

瑞士国家银行在 1974 年底定下 1975 年 M1 增长 6\% 的目标，想逐步而不是用休克疗法压下通胀\cite{rich2003}。1975 年，石油冲击加上紧缩的货币政策，使产出下降 7.5\%，是二战以后最急的下降\cite{rich2003}。石油冲击是\chref{ch:oil}的题目。
% 据: rich_1975, rich_gradual, rich_drop

1978 年，瑞士法郎升到前所未见的高度，出口部门、联邦政府与公众都向央行施加巨大的压力\cite{rich2003}。10 月初，瑞士国家银行宣布临时的汇率目标，要让德国马克（即联邦德国的货币）的瑞士法郎价格“明显”高于 80\cite{rich2003}。干预使基础货币与 M1 爆发式扩张，1978 年的目标超出 10 个百分点以上\cite{rich2003}。汇率目标的来龙去脉在\chref{ch:floating}。1988 年，瑞士的银行与央行推出世界上第一个实时全额结算系统，银行大大减少了持有的基础货币\cite{rich2003}。按基础货币定的目标因此失准，表中 1988—1990 年的负增长就是这样来的（本书的解读）。
% 据: rich_never, rich_pressure, rich_80, rich_explosive, rich_1988, rich_sic
\end{case}

两家央行怎样评价自己的目标？Rich 曾是瑞士国家银行经济部的主管\cite{rich2003}。他的结论是，瑞士国家银行靠货币目标把通胀的趋势降到了低水平\cite{rich2003}。可是在经济扩张时，目标没有要求足够先发制人的紧缩，物价稳定就保持得不够好\cite{rich2003}。1999 年底，瑞士国家银行放弃货币目标，改用以通胀预测为中心的办法\cite{rich2003}。联邦银行的年度目标公布到 1998 年为止\cite{gerberding2004}。
% 据: rich_title, rich_wp, rich_formerly, rich_trend, rich_preemptive, rich_1999, gwz_t1998

\section{争论}\label{sec:monetarism-debate}

\subsubsection{货币是不是收入的原因}

弗里德曼与施瓦茨给对手出了两道题\cite{friedmanschwartz1963}。任何别的解释，都要说明收入的大变动为什么发生在那些时候，也要说明它们为什么总伴随着货币增长率的相应变动\cite{friedmanschwartz1963}。他们认为货币解释同时答了这两道题\cite{friedmanschwartz1963}。单是第一道题，别的解释也很难答，除非说不同的因素不时经由不同的渠道造成收入的大变动，他们说那等于承认一无所知\cite{friedmanschwartz1963}。货币存量的重大变动来源五花八门，许多既与同期的商业状况无关，也与更早的商业状况无关\cite{friedmanschwartz1963}。在他们看来，这就驳倒了从收入到货币的单向解释\cite{friedmanschwartz1963}。
% 据: fs_tests, fs_ignorance, fs_variety, fs_contradicts

他们也承认反方向的渠道。美国的银行体系长期容易遭挤兑，这一点几乎就能解释，收入的大幅下降不管怎样产生，为什么通常伴随着货币的大幅下降\cite{friedmanschwartz1963}。可是他们说，这个解释并不适用于所有的下降，对货币的上升也基本不相干\cite{friedmanschwartz1963}。他们还反问，货币若只是被动的伴随者，中国怎么躲过了大萧条的头几年？他们的回答是，中国用银本位，对金本位国家的汇率是浮动的\cite{friedmanschwartz1963}。\chref{ch:depressionworld}讲过，对中国的这段经历另有不同的看法（\secref{sec:depressionworld-silver}）。
% 据: fs_runs, fs_notall, fs_china

同一本文集里，明斯基（Hyman P. Minsky，1919—1996）的评论不同意\cite{friedmanschwartz1963,loc45,gnd45}。他说，1870—1960 年记录在案的 21 个周期里有 6 个是深度萧条\cite{friedmanschwartz1963}。其中 1919—1920 年与 1937—1938 年的两次货币下降，可以认定是联邦储备的具体行动造成的，都没有伴随金融危机\cite{friedmanschwartz1963}。另外四次，货币的减少本身就是金融危机的一部分\cite{friedmanschwartz1963}。他认为，要把货币解释补完整，就得让金融危机成为体系内生的，而不是外来的\cite{friedmanschwartz1963}。至今没有定论。
% 据: mk_cycles, mk_administered, mk_crisis, mk_comment, mk_endogenous, loc_minsky, loc_minsky_head, gnd_minsky_b, gnd_minsky_d

\subsubsection{目标失败，是货币需求变了，还是目标没有守}

一方的理由来自央行自己的记录。联邦公开市场委员会、Bouey 与劳森都把失灵归于金融创新与管制的放松（\secref{sec:monetarism-goodhart}）。Bouey 在 1982 年 12 月 8 日对参议院的委员会说，加拿大银行放弃目标纯粹出于技术上的理由，并没有改变对货币目标在原则上可取的看法\cite{bouey1982b}。他还说，M1 目标从来没有被当作自动驾驶仪\cite{bouey1982b}。
% 据: boud_date, boud_technical, boud_pilot

另一方的理由是目标没有真正守过。Clarida 与 Gertler 等人质疑货币目标对联邦银行日常决策的意义\cite{gerberding2004}。Clarida、Galí 与 Gertler 把货币与目标路径的差距加进联邦银行的反应函数，发现货币总量“根本无关紧要”\cite{gerberding2004}。Gerberding、Worms 与 Seitz 用当时实际能看到的数据重估，发现货币增长偏离目标的程度也高度显著，联邦银行是认真对待目标的\cite{gerberding2004}。瑞士国家银行一方面公布目标，一方面不愿公开承认自己也看经济周期\cite{rich2003}。它 1980 年的基础货币目标是增长 4\%，起点定在 1979 年 11 月，这样就不必宣布一个负的目标\cite{rich2003}。新闻稿没有说明这个起点，Rich 说这是新闻稿漏掉关键信息的极少数例子之一\cite{rich2003}。至今没有定论。
% 据: gwz_cgg_challenge, gwz_cgg_test, gwz_cgg, gwz_realtime, rich_cycle, rich_1980, rich_negative, rich_omitted, rich_omitted2

\subsubsection{古德哈特定律会不会让任何中间目标落空}

Moses 1983 年设想，若改以信用总量为操作目标，古德哈特定律可能又会起作用\cite{moses1983}。利-彭伯顿 1986 年却说，对广义货币的表现重视不够，与对它强调得太多、太精确一样不明智\cite{leighpemberton1986}。劳森 1987 年不再为广义货币设目标，仍说会继续考虑它\cite{hansard1987}。瑞士国家银行改用通胀预测以后，M3 等货币总量仍是重要的指标\cite{rich2003}。至今没有定论。
% 据: moses_credit, lp86_unwise, lp86_unwise2, law87_account, rich_m3

\begin{deeper}[史料说明]
豪、劳森 1986 年与 1987 年的预算演说据英国议会的历史议事录网站；网页的最后一行带有每次不同的脚本，散列值是本书读的那一份。联邦公开市场委员会 1979 年 10 月 6 日、1982 年 2 月 1—2 日与 10 月 5 日的政策记录据联邦储备网站的 PDF，记录里写作“Ml”的是 M1。英格兰银行 1981 年 3 月的季报与利-彭伯顿的演讲据英格兰银行网站；季报的文字层把分数读错了，如 17.5\% 读成“17!-\%”，相关的数按页面图像核对。弗里德曼与施瓦茨的长文与明斯基的评论据美国国民经济研究局 1965 年出版的文集的扫描本，部分页面按页面图像核对。文中说两人的《美国货币史》正在付印，注明 1963 年，本书因此把这篇文章记作 1963 年。

古德哈特 1975 年的论文没有读，他的法则按 Moses（1983）的转引；后来流行的另一种更简短的说法另有出处，本书不用。普尔 1970 年的论文也没有读，按普尔与 Lieberman（1972）的复述。Rich 的表 1 在印刷页第 18 页，Gerberding 等人的表 10 在印刷页第 45 页。《总统经济报告》表 B-67 的 M1 是 12 月的日平均数，流通速度因此用 12 月的 M1 计算，与用全年平均的 M1 算出的数略有不同。联邦银行为什么在 1998 年以后不再公布目标，所读的文献没有说明。
% 据: hansard_copyright, fomc82f_ml, qb81_title, qb81_garble, fs_cover, fs_pubdate, moses_ref, erp_b67_dec

Bouey 的两篇讲话据加拿大银行的开放档案，来源是加拿大银行，引文由本书译成中文。Moses 的讨论稿据澳大利亚储备银行网站，按 CC BY 4.0 使用，来源：Reserve Bank of Australia 1983。Gerberding 等人的讨论稿据联邦银行网站，Rich 的工作论文据欧洲中央银行网站，本书在正文与表下都注明了出处。英格兰银行的材料只用数字与短语，不复制图表。弗里德曼 1968 年的文章与布鲁金斯学会的三篇文章也只引短语。
% 据: boc_permit, boc_attrib, rba_licence, rba_attrib, gwz_title, gwz_series, bundesbank_repro, ecb_copy, ecb_repro, boe_legal_copyright, boe_legal_permission, aea_terms, brookings_terms
\end{deeper}

\section{今天还剩什么}\label{sec:monetarism-legacy}

货币需求稳定这一条，1970 年代中期以后在美国、英国与加拿大都没有守住。长期中性那一条留了下来，成了\chref{ch:naturalrate}的长期竖线。单一规则的精神也留了下来，央行要公开宣布自己的目标，并为结果负责（本书的概括）。瑞士国家银行认为，通胀预测比它的中期货币战略更容易向公众解释\cite{rich2003}。
% 据: rich_communicate

\chref{ch:silver}留下了数量论的两个老问题，一是流通速度是否稳定，二是货币与物价谁先谁后。货币主义的二十年给第一个问题一个分开的回答，短期里流通速度不稳定，长期的趋势还可以估计（本书的解读）。第二个问题，弗里德曼、施瓦茨与明斯基的争论还没有了结。\chref{ch:bullion}讲过，在可兑换的金本位下货币是物价的结果，在不兑现的纸币下它可以是原因。1970 年代各国央行管的，正是不兑现的纸币（本书的概括）。

货币总量本身也没有消失，只是从目标降成了看的东西之一（本书的概括）。Moses 1983 年主张看一系列金融变量\cite{moses1983}。劳森 1987 年按狭义货币、广义货币与汇率定利率\cite{hansard1987}。弗里德曼 1968 年要的是一个公布出来、人所共知的增长率，后来的央行公布的是通胀的目标（本书的解读）。取代货币目标的通胀目标制与泰勒规则，在\chref{ch:consensus}。
% 据: moses_range, law87_rates, f68_known

\begin{takeaways}
\begin{itemize}
\item 货币主义的三个命题：货币需求稳定，货币按公开的规则稳定增长，货币长期中性。弗里德曼与施瓦茨说，货币存量的六次重大下降都对应着严重的经济收缩。
\item 货币目标由数量方程倒推，$m^{*}=\pi^{*}+y^{*}-v^{*}$。流通速度一旦不可预测，目标就失去了根据。
\item 普尔的规则：实体扰动为主时盯货币，货币扰动为主时盯利率。1970 年代中期以后美国的货币需求移动了，Goldfeld 据此主张利率政策。
\item 美国 1979 年 10 月改为控制准备金，1982 年 10 月不再为 M1 定具体目标。英国 1980 年的£M3 目标一再超出，1987 年放弃。加拿大 1982 年 11 月放弃 M1 目标。
\item 联邦银行与瑞士国家银行守着货币目标二十多年，目标常常落空，通胀的趋势却降了下来。古德哈特定律提醒，被当作目标的统计规律可能因此瓦解。
\end{itemize}
\end{takeaways}

\begin{selfcheck}
\begin{enumerate}
\item 能容忍的通胀是 2\%，潜在产出增长 3\%，流通速度每年下降 1\%，货币目标是多少？
\item 上题的目标照常执行，流通速度实际每年下降 3\%，通胀是多少？
\item 按\exref{ex:monetarism-poole}的参数，货币扰动的方差是 3 时，哪种工具让产出更稳？
\item 表 B-1 中 1981 年的名义 GDP 是 31313 亿美元，表 B-67 中 1981 年 12 月的 M1 是 4362 亿美元，流通速度是多少？
\item 瑞士国家银行 1978 年的 M1 目标是 5\%，实际增长 16.2\%，超出多少个百分点？
\end{enumerate}
\end{selfcheck}

\begin{answers}
\item $2+3-(-1)=6\%$。
\item $6+(-3)-3=0$，物价不涨也不跌。
\item $\sigma^{2}_{Y,M}=(0.25+3)/2.25\approx1.44$，大于 1，盯利率更好。
\item $31313/4362\approx7.18$。
\item 11.2 个百分点，与 Rich 说的“超出 10 个百分点以上”相符。
\end{answers}

回到 1980 年 3 月 26 日的下院。豪说，别的国家碰到同样的问题，也用了同样的办法，两个星期前美国就采取了坚决的措施\cite{hansard1980}。各国能各自管住本国的货币，前提是汇率不再固定。瑞士 1973 年让法郎浮动以后，央行才在原则上完全控制了货币供给\cite{rich2003}。1978 年法郎涨得太猛，它又不得不回头盯汇率\cite{rich2003}。货币目标的二十年，正好落在固定汇率瓦解之后（本书的解读）。美元与黄金为什么脱钩？浮动汇率为什么会“超调”？
% 据: h80_us, rich_floatdate, rich_float, rich_80
